\subsubsection{Segment Tree 2D}
\begin{lstlisting}[style=mystyle, language=C++]
// TC: O(log N * log M) per query
// usa: matriz con queries suma en subrectangulo y updates punto
#include <bits/stdc++.h>
using namespace std;

struct SegTree2D{
	int n, m;
	vector<vector<int>> tree;
	SegTree2D(int n, int m): n(n), m(m){
		tree.assign(4*n, vector<int>(4*m, 0));
	}
	void updateY(int vx, int lx, int rx, int vy, int ly, int ry, int x, int y, int val){
		if(ly == ry){
			if(lx == rx) tree[vx][vy] = val;
			else tree[vx][vy] = tree[vx*2][vy] + tree[vx*2+1][vy];
		}else{
			int my = (ly + ry) / 2;
			if(y <= my) updateY(vx, lx, rx, vy*2, ly, my, x, y, val);
			else updateY(vx, lx, rx, vy*2+1, my+1, ry, x, y, val);
			tree[vx][vy] = tree[vx][vy*2] + tree[vx][vy*2+1];
		}
	}
	void update(int x, int y, int val){
		updateY(1, 0, n-1, 1, 0, m-1, x, y, val);
	}
	int queryY(int vx, int vy, int ly, int ry, int y1, int y2){
		if(y1 > ry || y2 < ly) return 0;
		if(y1 <= ly && ry <= y2) return tree[vx][vy];
		int my = (ly + ry) / 2;
		return queryY(vx, vy*2, ly, my, y1, y2) + queryY(vx, vy*2+1, my+1, ry, y1, y2);
	}
	int queryX(int vx, int lx, int rx, int x1, int x2, int y1, int y2){
		if(x1 > rx || x2 < lx) return 0;
		if(x1 <= lx && rx <= x2) return queryY(vx, 1, 0, m-1, y1, y2);
		int mx = (lx + rx) / 2;
		return queryX(vx*2, lx, mx, x1, x2, y1, y2) + queryX(vx*2+1, mx+1, rx, x1, x2, y1, y2);
	}
	int query(int x1, int y1, int x2, int y2){
		return queryX(1, 0, n-1, x1, x2, y1, y2);
	}
};

int main(){
	SegTree2D st(8, 8);
	st.update(2, 3, 5);
	st.update(5, 5, 7);
	cout << st.query(0, 0, 7, 7) << "\n";  // 12
	return 0;
}
\end{lstlisting}
